\documentclass[11pt]{article}
\usepackage[margin=1.05in]{geometry}
\usepackage{amsmath,amssymb,amsthm}
\usepackage{booktabs}
\usepackage{enumitem}
\usepackage{xcolor}
\definecolor{linkblue}{RGB}{25,60,120}
\usepackage[colorlinks=true,linkcolor=linkblue,citecolor=linkblue,urlcolor=linkblue]{hyperref}
\usepackage[capitalize]{cleveref}
\emergencystretch=2em

\newtheorem{theorem}{Theorem}[section]
\newtheorem{lemma}[theorem]{Lemma}
\newtheorem{proposition}[theorem]{Proposition}
\newtheorem{corollary}[theorem]{Corollary}
\theoremstyle{definition}
\newtheorem{definition}[theorem]{Definition}
\newtheorem{problem}[theorem]{Problem}
\theoremstyle{remark}
\newtheorem{remark}[theorem]{Remark}

\newcommand{\F}{\mathbb F}
\newcommand{\RS}{\mathrm{RS}}
\newcommand{\emca}{\varepsilon_{\mathrm{mca}}}
\newcommand{\eca}{\varepsilon_{\mathrm{ca}}}
\newcommand{\Lst}{\Lambda}
\newcommand{\Ht}{\mathrm H_2}
\newcommand{\Paid}{\operatorname{Paid}}
\newcommand{\ceil}[1]{\left\lceil #1\right\rceil}
\newcommand{\floor}[1]{\left\lfloor #1\right\rfloor}

\title{Threshold Certificate Compilers for the CAP25 Reed--Solomon Program}
\author{RS-MCA experimental note\\ \small integrated from AllenGrahamHart PRs, with Codex editing}
\date{July 3, 2026}

\begin{document}
\maketitle

\begin{abstract}
This note packages the threshold and compiler material that is useful for a
future revision of the CAP25/Paper D threshold paper.  It is deliberately not a
new MCA theorem.  Its role is to make row-specific certificates mechanically
auditable: numerator staircases, safe/unsafe corridors, already-paid tangent,
quotient, and extension cells, dyadic quotient census windows, dodge-selection
arithmetic, and integrality margins are all stated as typed compiler lemmas.

The main output is a reusable certificate grammar.  Given a row-specific packet
with a denominator, certified lower and upper numerators, and branch coverage,
the grammar tells where the first safe agreement may lie, which already-paid
terms may be charged, which quotient census windows remain unresolved, and when
large entropy margins make polynomial residual factors irrelevant.  The note is
self-contained at the level of these arithmetic compilers; it cites the
Reed--Solomon and proximity background papers only for the code-theoretic
objects and the already-proved ledgers.
\end{abstract}

\section{Purpose and Status}

Paper D/CAP25 reduces a prize-facing threshold claim to exact integer ledgers:
one prints a row, a radius convention, a denominator, an unsafe lower numerator,
and a safe upper numerator or residual certificate.  The threshold work in this
note provides the arithmetic layer around those row-specific packets.  It should
be read as \emph{proved compiler arithmetic} plus a few experimental margin
tables, not as a proof of the missing aperiodic safe-side theorem.

\begin{remark}[How this fits Paper D v12]
The note supplies supporting infrastructure for the certificate grammar of
Paper D v12.  It should be used in a future CAP25 revision to replace informal
phrases such as ``the crossing is adjacent'' or ``the quotient window is small''
by explicit lemmas.  It does not alter Paper D's universal cap theorem, its
safe-side reductions, or any deployed-row theorem by itself.
\end{remark}

\section{Integer Staircases}

\begin{definition}[integer numerator staircase]
Fix an agreement interval
\[
        I=\{a_{\min},a_{\min}+1,\ldots,a_{\max}\}
\]
and a nonincreasing function
\[
        N:I\longrightarrow \mathbb Z_{\ge 0}.
\]
For MCA or line-MCA, $N(a)$ is the number of bad sampled parameters at
agreement at least $a$; the closed Hamming radius is $r=n-a$.  For list
decoding, one may use agreement coordinates $a=n-r$, so the list numerator is
also nonincreasing in $a$.

Given a denominator $Q$ and target $\varepsilon^*$, put
\[
        B_*=\floor{\varepsilon^* Q}.
\]
An agreement $a$ is \emph{safe} if $N(a)\le B_*$ and \emph{unsafe} if
$N(a)>B_*$.
\end{definition}

\begin{theorem}[staircase localization]\label{thm:staircase}
Let $N:I\to\mathbb Z_{\ge0}$ be nonincreasing.  Exactly one of the following
holds:
\begin{enumerate}[label=(\roman*)]
        \item all $a\in I$ are safe;
        \item all $a\in I$ are unsafe;
        \item there is a unique first safe agreement $a_*\in I$ such that
        \[
                N(a_*-1)>B_*\ge N(a_*),
        \]
        with $a_*-1\in I$.
\end{enumerate}
In the interior case the safe set is the upper interval $\{a\ge a_*\}$ and the
unsafe set is the lower interval $\{a<a_*\}$.
\end{theorem}

\begin{proof}
The predicate $N(a)\le B_*$ is monotone nondecreasing in $a$.  A monotone
Boolean function on a finite interval is identically false, identically true,
or has a unique first true point.
\end{proof}

\begin{corollary}[closed-radius endpoint convention]\label{cor:endpoint}
In the interior MCA case, the largest safe closed integer radius is
\[
        r_{\mathrm{safe}}=n-a_*,
\]
and the first unsafe closed integer radius is $r_{\mathrm{safe}}+1$.  As a
real closed-ball threshold, the supremum of safe radii is
\[
        \frac{n-a_*+1}{n},
\]
but that endpoint is not attained if $N(a_*-1)>B_*$.
\end{corollary}

\begin{theorem}[corridor from lower and upper certificates]\label{thm:corridor}
Suppose $L(a)\le N(a)\le U(a)$ are certified lower and upper bounds on the same
staircase.  Define
\[
        A_{\mathrm{unsafe}}=\{a\in I: L(a)>B_*\},
        \qquad
        A_{\mathrm{safe}}=\{a\in I: U(a)\le B_*\}.
\]
Every interior first safe agreement satisfies
\[
        1+\max A_{\mathrm{unsafe}}\le a_* \le \min A_{\mathrm{safe}},
\]
with the conventions $1+\max\varnothing=a_{\min}$ and
$\min\varnothing=a_{\max}$.
\end{theorem}

\begin{proof}
If $L(a)>B_*$ then $N(a)>B_*$, so the first safe point is strictly after $a$.
If $U(a)\le B_*$ then $N(a)\le B_*$, so the first safe point is at most $a$.
Taking the strongest certified unsafe and safe agreements gives the interval.
\end{proof}

\begin{theorem}[coarse steepness compiler]\label{thm:steepness}
Let $M(a)>0$ be a positive model count and suppose a row packet proves
\[
        \frac{M(a)}{E}\le N(a)\le E M(a)
        \qquad(E\ge1)
\]
on consecutive agreements.  If
\[
        \frac{M(a)}{M(a+1)}>E^2,
\]
then the ambiguity intervals $[M(a)/E, E M(a)]$ and
$[M(a+1)/E, E M(a+1)]$ are disjoint.  Hence a fixed budget $B_*$ can lie in at
most one consecutive ambiguity interval, and all other agreements have their
safe/unsafe sign certified by the coarse envelope.
\end{theorem}

\begin{proof}
The separation hypothesis is exactly $M(a)/E>E M(a+1)$.  The sign implications
are immediate from the lower and upper envelope.
\end{proof}

\section{Paid Ledger Functions}

The CAP25 certificate grammar should separate branches that are already paid
from residual branches that require new structural input.

\subsection{Tangent cell}

\begin{definition}[high-agreement tangent cell]
For a Reed--Solomon row of length $n$ and dimension $k$, write $r(A)=n-A$ and
\[
        R_{\tan}=\floor{\frac{n-k}{3}}.
\]
Define the partial paid tangent cell
\[
        \Paid_{\tan}^{\mathrm{hi}}(n,k,A)=r(A)+1
\]
when $0\le r(A)\le R_{\tan}$, and leave it unavailable otherwise.
\end{definition}

\begin{proposition}[tangent cell usage]\label{prop:tangent}
In the high-agreement range $0\le n-A\le\floor{(n-k)/3}$, the finite affine
support-wise MCA line numerator is exactly $n-A+1$.  Therefore
$\Paid_{\tan}^{\mathrm{hi}}$ is an exact paid cell in this range.  Outside this
range it should not be used without a separate tangent/common-line theorem.
\end{proposition}

\begin{remark}
For the row $n=512$, $k=256$, and denominator $Q=17^{32}$, one has
$\floor{Q/2^{128}}=6$, while
\[
        \Paid_{\tan}^{\mathrm{hi}}(512,256,506)=7,\qquad
        \Paid_{\tan}^{\mathrm{hi}}(512,256,507)=6.
\]
This recovers the high-agreement exact gate in the finite-slope convention.
\end{remark}

\subsection{Quotient cell}

Let $C$ be a finite set of quotient fiber sizes $c\mid n$.  A conservative
support upper cell at agreement $A$ is
\[
U_{\mathrm{sum}}(n,A,C)=
\sum_{c\in C}\sum_{B=A}^{n}
\binom{n/c}{\floor{B/c}}
\binom{n-c\floor{B/c}}{B-c\floor{B/c}}.
\]
If an exact support-union or image-lcm certificate is available, that sharper
integer may replace $U_{\mathrm{sum}}$.

\begin{definition}[quotient paid status]
The quotient paid cell at $(n,A,C)$ is status-valued:
\[
\Paid_{\mathrm{quot}}=
\begin{cases}
\mathrm{EXACT\_IMAGE}(\deg R) & \text{if an image lcm certificate is supplied},\\
\mathrm{EXACT\_SUPPORT}(U_{\mathrm{supp}}) & \text{if a block-profile union is supplied},\\
\mathrm{SAFE\_SUM}(U_{\mathrm{sum}}) & \text{otherwise}.
\end{cases}
\]
\end{definition}

\begin{remark}[zone tags]
Unsafe quotient lower cells should carry one of the following tags:
\[
\begin{array}{ll}
\mathrm{ZONE\_A\_NORM\_EXACT} & \text{the norm-exact gate holds},\\
\mathrm{ZONE\_B\_COLLISION\_OPEN} & \text{collisions remain unresolved},\\
\mathrm{ZONE\_C\_CAP\_FLOOR} & \text{a cap/floor theorem supplies the mass}.
\end{array}
\]
Zone B is interval-valued.  A CAP25 table should never silently replace it by a
point estimate.
\end{remark}

\subsection{Extension cell}

Let $B$ be the generated field and $F$ the line field, with
$|B|=q_{\mathrm{gen}}$ and $|F|=q_{\mathrm{line}}$.

\begin{proposition}[extension paid cell]\label{prop:extension}
If $q_{\mathrm{gen}}=q_{\mathrm{line}}$, then there is no proper extension-only
line-parameter cell:
\[
        \Paid_{\mathrm{ext}}^{\mathrm{only}}=0.
\]
If $q_{\mathrm{gen}}<q_{\mathrm{line}}$, the extension-pole lower mechanism
maps a base list lower bound $L$ to
\[
        \mathrm{ExtPole}(q_{\mathrm{line}},q_{\mathrm{gen}},k,L)
        =
        \left\lceil
        \frac{L(q_{\mathrm{line}}-q_{\mathrm{gen}})}
        {q_{\mathrm{line}}-q_{\mathrm{gen}}+kL}
        \right\rceil .
\]
For safe-side extension charts over $B$, a closed parameter set of degree
$\Delta$ and dimension $e$ contributes at most $\Delta q_{\mathrm{gen}}^e$
affine extension parameters.
\end{proposition}

\section{Dyadic Quotient Census}

For a $2$-power quotient order $N'=2^a$, let $n_1=N'/2$.  The characteristic
zero antipodal quotient count used by Paper B is
\[
        A_2(N',\ell')=
        \sum_{\substack{u\ge0\\ t=\ell'-2u\ge0\\u\le n_1-t}}
        \binom{n_1}{t}2^t.
\]
At rate $1/2$, where $\ell'=n_1+1$,
\[
        A_2(N',n_1+1)=\frac{3^{n_1}-1}{2}.
\]

\begin{theorem}[bounded quotient census]\label{thm:census}
Let $B_{\max}=2^{128}-1$, corresponding to the largest possible
$\floor{q/2^{128}}$ under $q<2^{256}$.  In the relaxed even-scale model, the
first scales at which $A_2(N',\rho N'+1)>B_{\max}$ are
\[
\begin{array}{c|cccc}
\rho & 1/2 & 1/4 & 1/8 & 1/16\\ \hline
N'_{\mathrm{relaxed}} & 164 & 176 & 248 & 384 .
\end{array}
\]
For actual dyadic quotient orders, the corresponding first dyadic scales are
\[
\begin{array}{c|cccc}
\rho & 1/2 & 1/4 & 1/8 & 1/16\\ \hline
N'_{\mathrm{dyadic}} & 256 & 256 & 256 & 512 .
\end{array}
\]
Thus the exact integer quotient endgame sees only bounded quotient scales,
independent of the deployed blocklength, once the row has been reduced to this
canonical quotient-count model.
\end{theorem}

\begin{remark}
The theorem is a compiler fact, not a collision theorem.  It says which finite
scales must be inspected; it does not prove that a row-specific lower packet
survives reduction modulo a finite field.
\end{remark}

\begin{theorem}[budget windows and dodge selection]\label{thm:windows}
Let $K$ be an exact count and $L\le K$ a certified lower count for the same
object.  Put $B(q)=\floor{q/2^{128}}$.  Then
\[
        B(q)<L \Rightarrow \text{certified unsafe},\qquad
        B(q)\ge K \Rightarrow \text{certified safe},
\]
and the only unresolved budgets are
\[
        L\le B(q)<K.
\]
Equivalently the unresolved field-size interval is
\[
        L2^{128}\le q\le K2^{128}-1.
\]
If the lower proof is $L=K-g$, then the largest tolerable missing mass while
still proving unsafe is
\[
        g_{\max}=K-B(q)-1.
\]
\end{theorem}

\begin{proof}
All assertions are restatements of the integer inequalities $L\le N\le K$ and
the target condition $N\le \floor{q/2^{128}}$.
\end{proof}

\section{Dyadic Quotient Profile}

For $n=2^m$, $k=\rho n$, and integer slack
$\sigma=\ceil{\eta n}$, define the exact dyadic quotient-profile count
\[
K_Q(n,k,\sigma)=
\max_{\substack{M=2^e\mid \gcd(n,k)\\ \sigma<M\\ k/M\le n/M-1}}
\binom{n/M-1}{k/M},
\]
with value $0$ if the active set is empty.  This is the integer form of the
Paper B quotient profile and avoids floating-point logarithms.

\begin{proposition}[bounded active quotient order]\label{prop:qprofile}
If $\sigma\ge n/256$, then every active scale in $K_Q$ has quotient order
$N=n/M\le128$.  If $\sigma\ge n/128$, then every active scale has
$N\le64$.
\end{proposition}

\begin{proof}
Activity requires $\sigma<M=n/N$.  Therefore $N<n/\sigma$.  The displayed
claims follow from the two lower bounds on $\sigma$.
\end{proof}

\section{Integrality Margins}

The integrality-margin tables are experimental arithmetic evidence, but their
logical use is simple and worth isolating.

\begin{proposition}[polynomial residual absorption]\label{prop:integrality}
Fix a candidate scale and suppose a future theorem bounds an unpaid residual
contribution by
\[
        n^C \Phi(n,q)
\]
for some explicit proxy $\Phi$.  If the verified entropy calculation gives
\[
        -\log_2(n^C\Phi(n,q))>0,
\]
then the residual integer contribution is zero.
\end{proposition}

\begin{proof}
The residual contribution is a nonnegative integer bounded above by a real
number strictly smaller than $1$.
\end{proof}

\begin{remark}[recorded margin]
The integrated verifier evaluates max official-dimension rows
$k=2^{40}$, $n\in\{2k,4k,8k,16k\}$, and budget bits $64,96,128$.  For the
tested candidate crossing scales, even an $n^3$ multiplier is absorbed with
very large margins; the worst recorded margin is about $9.68\cdot10^7$ bits.
This says the hard remaining question is structural classification, not
constant sharpening.
\end{remark}

\section{CAP25 Packet Schema}

A future CAP25 threshold packet can cite this note by printing the following
fields:
\[
\begin{array}{ll}
\text{row} & (F,D,k,n,\rho),\\
\text{denominators} & q_{\mathrm{gen}},q_{\mathrm{line}},q_{\mathrm{chal}},\\
\text{target} & \varepsilon^*,\quad B_*=\floor{\varepsilon^* Q},\\
\text{agreement interval} & I=[a_{\min},a_{\max}]\cap\mathbb Z,\\
\text{unsafe certificates} & L(a)>B_*,\\
\text{safe certificates} & U(a)\le B_*,\\
\text{paid cells} & \Paid_{\tan},\Paid_{\mathrm{quot}},\Paid_{\mathrm{ext}},\\
\text{residual cells} & \text{named aperiodic/extension/list branch},\\
\text{deduplication rule} & \text{support/image/root coalescing theorem},\\
\text{endpoint convention} & \text{closed integer ball and real supremum}.
\end{array}
\]
Then \Cref{thm:staircase,thm:corridor} locate the remaining threshold
corridor; \Cref{thm:windows} translates quotient or planted counts into exact
field-size windows; and \Cref{prop:integrality} turns sufficiently strong
entropy margins into zero residual integer contributions.

\section{Non-Claims}

This note does not prove an aperiodic M1 local limit, an L1 image-fiber upper
bound, an extension-line transfer theorem, a protocol soundness theorem, or a
new Paper D cap.  It is a compiler layer: once those mathematical inputs are
available, it makes the threshold conclusion exact and auditable.

\begin{thebibliography}{9}
\bibitem{ABF26}
S. Arnon, D. Boneh, and N. Fenzi.
\newblock Open problems in list decoding and correlated agreement.
\newblock ePrint 2026/680, 2026.

\bibitem{Cho26B}
P. Chojecki.
\newblock Smooth-domain Reed--Solomon slack MCA and quotient-profile ledgers.
\newblock RS-MCA repository, Paper B, 2026.

\bibitem{Cho26D}
P. Chojecki.
\newblock Universal field-size caps and a two-sided sandwich for mutual
correlated agreement on smooth Reed--Solomon domains.
\newblock RS-MCA repository, Paper D/CAP25 v12, 2026.
\end{thebibliography}

\end{document}
